\chapter{理想链模型与基本统计量}
\section{模型定义}
\subsection{链段与末端矢量}
\subsubsection{独立链段假设}
设一条链由$N$个长度为$b$的独立链段组成，链段取向各向同性。末端矢量及其均方值分别为
\begin{align}
  \boldsymbol{R} &= \sum_{i=1}^{N}\boldsymbol{b}_i,\\
  \langle R^2\rangle &= \sum_{i=1}^{N}\langle b_i^2\rangle
  +2\sum_{i<j}\langle\boldsymbol{b}_i\cdot\boldsymbol{b}_j\rangle=Nb^2.
\end{align}
交叉项为零来自统计独立性与各向同性；对于真实链，应先检查局部相关性是否可通过有效链段吸收。

\subsubsection{回转半径}
长理想链满足$\langle R_g^2\rangle=Nb^2/6$。有限链长修正与链段定义有关，因此比较不同模型时必须保持统计量定义一致。

\section{模型的边界}
独立链段模型忽略了排除体积。讨论良溶剂中的链尺寸时，需要重新考虑构象之间的相互作用，不能只改变拟合参数而保持原有解释。\footnote{参见 de Gennes P G. Scaling Concepts in Polymer Physics[M]. Ithaca: Cornell University Press, 1979.}
本章与引言的注码连续递增，用于检查跨章脚注编号。
